\begin{textsample}{POS Dim 1 – al – Score 37.00 – al/al\_em\_40.txt}  \label{ex:f1_pos_155}
HISTÓRIA “ A incompreensão do presente nasce fatalmente da ignorância do passado ”.
Com estas palavras, o historiador Marc Bloch ( 2001, p. 65 ) expressou \textbf{uma} premissa paradoxal no campo da História : a de \textbf{que} a compreensão do presente não se esgota em si mesmo, no seu próprio tempo, mas \textbf{exige} o exame do passado.
Eis \textbf{uma} contribuição ética, social e pública da História.
A partir da análise do passado lançamos luzes para os dilemas e questões do presente.
E, por outro \textbf{lado}, é paradoxal na medida em \textbf{que} a explicação sobre o passado também é \textbf{atravessado} pelas questões do presente.
Assim \textbf{uma} das maiores contribuições \textbf{que} o ofício do historiador tem dado a sociedade brasileira \textbf{contemporânea}, é levar as pessoas à compreensão do presente a partir dos ensinamentos do passado.
Essa premissa permitiu \textbf{que} a história, \textbf{enquanto} \textbf{disciplina} escolar, promovesse o estudo/reflexão não do passado pelo passado, mas de diferentes temporalidades e espaços.
Aliás, o passado, como já alertou Marc Bloch, não é o objeto da História.
Seu objeto é o ser humano, ou melhor os seres humanos no tempo.
Ao mesmo tempo, a inclusão das experiências históricas dos agentes sociais nas discussões historiográficas tem feito com \textbf{que} o olhar da História se desloque para o estudo das sociedades, trajetórias humanas, e experiências socioculturais ao longo do tempo e espaço, \textbf{abordagens} no ensino e na escrita da História \textbf{que} também contemplem as histórias, culturas e sociedades não europeias, não ocidentais.
Emergem como temas de estudo, os grupos subalternos, as histórias locais, a diversidade cultural, o cotidiano, as mentalidades e as sensibilidades, entre outros tantos aspectos \textbf{que} contemplem a pluralidade de formas do fazer humano e social.
Vale destacar \textbf{que} o professor e a professora devem promover meios para \textbf{que} saberes outros sejam considerados e \textbf{que} os “ outros ” não sejam tomados como \textbf{simples} objetos de estudos, mas como agentes sociais \textbf{que} devem ser respeitados e valorizados.
A história como \textbf{disciplina} escolar no Brasil \textbf{remete} ao \textbf{século} XIX com a fundação do Colégio Pedro II em 1837.
Este promovia \textbf{um} ensino de História Geral organizado na perspectiva positivista, evolutiva e eurocêntrica.
Posteriormente incluiu -se História da Pátria ou do Brasil.
Durante o período imperial, a História Sagrada era parte integrante do programa de ensino ( BITTENCOURT : 2008 : 78 ).
Portanto era \textbf{um} ensino memorialista, de \textbf{forte} tradição europeia, ensinado com \textbf{uma} didática pautada em \textbf{métodos} de \textbf{memorização}, e \textbf{que} em grande medida perdurou até o início do \textbf{século} XX.
Se em muitos momentos a História esteve a serviço dos grandes heróis nacionais, dos chefes políticos, da construção da ideia de nação e caráter de patriotismo ( GUIMARÃES : 2010 ), como foi o caso do início da Primeira República, episódios como o Golpe de 1964 e a ditadura civil-militar \textbf{que} se seguiu no Brasil por 21 anos desafiou a História a resistir à moral e à cívica impostos pelo então poder \textbf{institucionalizado}.
A redemocratização trouxe consigo a necessidade de afirmação das liberdades, da democracia e da cidadania.
A História portanto, é chamada a exercer sua função crítica na consolidação da democracia no Brasil.
Em resposta a essas necessidades e às transformações político-sociais, em 1996 é promulgada a Lei de Diretrizes e Bases da Educação \textbf{que} em seu artigo 3º, define \textbf{que} o ensino escolar no Brasil terá como base alguns princípios como o “ respeito à liberdade e apreço à tolerância ” e “ vinculação entre a educação escolar, o trabalho e as práticas sociais ”. \textbf{Aqui} a História se insere como aporte para concretizar a ideia de liberdade e de tolerância às múltiplas culturas e povos \textbf{que} compõem o território brasileiro, assim como refletir as práticas sociais através das relações de trabalho.
Daí a necessidade de \textbf{que} o conhecimento escolar ultrapasse as barreiras teórico-conceituais.
Estudar o passado pelo passado, memorizar e reproduzir os eventos históricos já não satisfaziam e não satisfazem mais essa nova sociedade e suas demandas.
Consolidada a democracia no Brasil e frente aos desafios da educação de \textbf{hoje}, o ensino de História em Alagoas deve buscar desenvolver nos/as estudantes a capacidade de “ propor e questionar hipóteses sobre as ações dos agentes sociais e também, identificar ambiguidades e contradições presentes tanto nas condutas individuais como nos processos e estruturas sociais ” ( BRASIL : 2018 : 548 ).
Para atingir esse objetivo, os professores e as professoras de História do ensino médio têm o desafio de levar os/as estudantes a desenvolverem habilidades \textbf{que} os tornem capazes de ter o domínio de conceitos e metodologias próprias da ciência histórica, assim como o reconhecimento das possibilidades de atualizações de experiências históricas do passado no presente.
Enseja -se construções de alternativas de possibilidades de futuro e cidadania.
E, portanto, “ aprender a indagar, ponto de partida para \textbf{uma} reflexão crítica, é \textbf{uma} das contribuições essenciais das Ciências Humanas e Sociais Aplicadas ” e \textbf{um} dos princípios fundamentais da ciência histórica, \textbf{que} contribui decisivamente “ para o desenvolvimento da autonomia dos \textbf{sujeitos} ” ( BRASIL : 2018 : 549 ).
Compreende -se \textbf{que} essa autonomia é o caminho \textbf{norteador} para o desenvolvimento de \textbf{um} cidadão comprometido e atuante em sua sociedade.
O ensino de História contempla como eixos \textbf{norteadores} os conceitos de aprendizado histórico e consciência histórica ( RÜSEN : 2012 ).
Estes se tornam significativos ao instrumentalizarmos os/as estudantes a construir o conhecimento por meio do confronto de diferentes ideias advindas da análise e compreensão das fontes históricas.
Para Barca ( 2018 ), o processo de instrumentalização passa pela compreensão contextualizada do passado, pela interpretação das fontes históricas no presente, visando \textbf{uma} perspectiva de ação temporal para o futuro, consequência da compreensão desse passado no presente.
A História pretende contribuir para \textbf{que} o/a estudante possa atuar de forma cidadã frente às questões sociais, políticas, econômicas, culturais, raciais, de gênero e ambientais \textbf{que} o Brasil e, mais especificamente, o Estado de Alagoas, apresentam na \textbf{contemporaneidade}.
Desta forma, entendemos \textbf{que} o conhecimento histórico segue \textbf{uma} dinâmica de vivacidade dos processos sociais e contribui para explicar a realidade, ao mesmo tempo \textbf{que} serve de instrumento de motivação e orientação para transformá -la, seja do ponto de vista coletivo ou individual do agente social, inserindo -se como cidadão pleno, agindo em seus direitos e deveres na construção/ transformação da sociedade.
O ensino de História em Alagoas deve, portanto, pautar -se em investigar aspectos da \textbf{territorialidade} e das ações humanas nesse espaço nos diversos períodos de tempo.
Dessa forma, o professor e a professora de História de ensino médio devem apropriar -se dos \textbf{métodos} e das \textbf{teorias} da história para levar os/as estudantes a compreender a função e a utilização desse conhecimento.
É \textbf{preciso} \textbf{que} estudantes entendam o lugar social e \textbf{amparado} pelos \textbf{métodos} de construção do conhecimento histórico, sugere -se \textbf{que} construam análises \textbf{que} possibilite reflexão sobre sua realidade histórico-social, de forma a sentir -se agente desse processo.
Propõe -se então \textbf{uma} história \textbf{vista} de baixo ( SHARPE : 1992 ), permitindo o exame e uso das produções historiográficas \textbf{que} se propõem a recuperar a história dos grupos heterogêneos e subalternos e assim estudar \textbf{uma} História mais próxima da realidade dos/as estudantes.
Esse outro olhar sobre a História de Alagoas deverá ser o ponto de partida para estudar a História do Brasil e do Mundo.
Esta perspectiva de trabalho abre espaço para abordar temporalidades múltiplas e simultâneas delineando referências históricas \textbf{que} contemplem a pluralidade étnica, cultural, social, econômica e política de Alagoas.
E do mesmo modo reconhecer os diferentes grupos e agentes tradicionalmente identificados na historiografia e no ensino como subalternos e excluídos.
Neste sentido, questões \textbf{que} se levantam, são por exemplo, aquelas relacionadas aos significados do trabalho em diferentes sociedades, suas especificidades e experiências históricas.
Ao analisar as relações de produção e cultura de classe, comparado circunstâncias históricas, pode -se discutir as formas de exploração e resistência de índios, negros, sertanejos, \textbf{ribeirinhos}, operários etc., e os seus impactos na construção e transformação da sociedade alagoana.
Abre -se desse modo caminho para : “ conhecer e combater as diversas formas de desigualdade e violência, adotando princípios éticos, democráticos, inclusivos e solidários, e respeitando os Direitos Humanos ” ( BRASIL : 2018 : 564 ).
Desse modo, os/as estudantes do ensino médio poderão assumir atitudes de enfrentamento dos dilemas da vida em sociedade, posicionando -se como protagonista nas transformações políticas, econômicas e culturais.
Tem -se assim \textbf{uma} preocupação no ensino em respeitar “ diferentes posições, com vistas a possibilitar escolhas alinhadas ao exercício da cidadania e ao seu projeto de vida, com liberdade, autonomia, consciência crítica e responsabilidade ” ( BRASIL : 2018 : 565 ).
SOCIOLOGIA A Sociologia não possui \textbf{um} paradigma explicativo dominante, mas sim diferentes modelos \textbf{que} se dispõem a analisar as transformações e a complexidade das relações sociais ( LEAL ; YUNG, 2015 ).
A Sociologia escolar, constituindo -se do ensino da Sociologia, da Ciência Política e da Antropologia, pauta -se na transposição didática dessa pluralidade \textbf{teórica} e \textbf{metodológica}, visando explicar os fenômenos sociais, políticos e culturais.
No Brasil, a Sociologia vigorou de forma intermitente nos currículos escolares desde sua primeira inclusão no ensino secundário, em 1891.
Em 1925, a Sociologia foi introduzida no curso secundário do Colégio Pedro II, no Rio de Janeiro, como \textbf{disciplina} obrigatória no 6º ano ginasial.
Em 1930 passou a ser obrigatória em todas as escolas secundárias do país, o \textbf{que} se manteve até 1942.
Seu retorno deu -se gradativamente a partir de 1984 nos estados brasileiros e, em 2008 retornou como \textbf{disciplina} obrigatória em todo o ensino médio do país.
Para ensinar Sociologia é necessário \textbf{que} o professor, ou a professora, mobilize saberes \textbf{que} conduzam o/a estudante a pensar “ sociologicamente ” ( BAUMAN ; MAY, 2010 ), e o faça superar, ainda \textbf{que} minimamente, com o senso comum, de modo \textbf{que} se possibilite o desvelamento social e promova a interpretação das relações sociais a partir das contribuições \textbf{teórico-metodológicas} específicas do campo das Ciências Sociais.
O ensino de Sociologia pode ter como ponto de partida \textbf{um} dos quatro princípios \textbf{metodológicos} de ensino, sendo eles : temas, conceitos, \textbf{teorias} e pesquisas. \textbf{Contudo}, o seu ensino \textbf{demanda} mobilização dos conceitos, das \textbf{teorias} e dos \textbf{métodos} sociológicos, independente do ponto de partida \textbf{que} o professor, ou a professora, vier adotar. \textbf{Contudo}, é \textbf{desejável} \textbf{que} o ponto de partida da aula seja a realidade dos/as estudantes e \textbf{que} a condução das atividades de ensino aprendizagem possibilitem reflexões sociológicas \textbf{que} possam superar o estágio inicial de conhecimento dos/as estudantes.
A Sociologia, e consequentemente, a Sociologia Escolar, pode dotar os educandos de competências básicas \textbf{que} os capacitem a romper com \textbf{uma} visão estática e dramática da vida social, promovendo \textbf{uma} vigilância intelectual crítica ( FERNANDES, 1976 ). \textbf{Uma} das potencialidades da \textbf{disciplina} de Sociologia é promover condições para \textbf{que} os/as estudantes tenham a compreensão das relações existentes entre os indivíduos e a sociedade, sendo essa \textbf{uma} competência fundamental para o entendimento da realidade social.
Desse modo, o \textbf{educando} é levado a se perceber dentro de \textbf{uma} realidade cultural, social e histórica particular, rompendo com perspectivas lineares/ evolucionistas da história.
Se faz necessário considerar as muitas possibilidades de \textbf{interdisciplinaridade} entre a Sociologia e as demais \textbf{disciplinas} \textbf{que} compõem o currículo do ensino médio. \textbf{Contudo}, mantendo -se os respectivos marcos epistemológicos de cada campo do saber ( LIMA ; COSTA ; PERNAMBUCO, 2012 ).
Nessa \textbf{direção}, \textbf{uma} possibilidade de \textbf{interdisciplinaridade} está em escolher temas comuns com outras áreas e desenvolver \textbf{abordagens} e pesquisas a partir de orientações \textbf{teórico-metodológicas} diversas, sendo tal esforço \textbf{um} dos caminhos importantes para o diálogo entre as \textbf{disciplinas} e \textbf{um} ensino-aprendizagem mais completo e significativo, já \textbf{que} a realidade social é multifacetada.
A Sociologia Escolar tem \textbf{um} papel importante na formação pessoal e social dos educandos, \textbf{uma} vez \textbf{que} compreender o mundo social possibilita a compreensão de si próprio, de sua trajetória e seu estar no mundo.
Seu contributo estar também em proporcionar condições para o entendimento do mundo social, do desenvolvimento do senso crítico e da capacidade empreendedora dos nossos/as estudantes frente aos desafios da sociedade brasileira e alagoana, bem como propiciar condições para \textbf{que} o \textbf{educando} estabeleça conexões entre a sua visão de mundo e a estrutura social, compreendendo de quais formas sua biografia está ligada às biografias de outros agentes ( BAUMAN ; MAY, 2010 ) e como as estruturas sociais se relacionam com suas práticas e condições sociais ( MILLS, 1969 ).

% matched lemmas: abordagem, amparar, aqui, atravessar, contemporaneidade, contemporâneo, contudo, demandar, desejável, direção, disciplina, educar, enquanto, exigir, forte, hoje, institucionalizar, interdisciplinaridade, lado, memorização, metodológico, método, norteador, preciso, que, remeter, ribeirinho, simples, sujeito, século, teoria, territorialidade, teórico, teórico-metodológico, um, ver
\end{textsample}
